% Bewertungspaket zum Gesamttest «Trigonometrische Gleichungen» — Quelle des PDFs
% (python3 scripts/build-lp-pdf.py trigonometrische-gleichungen). Ganze Punkte je Teilschritt; Werte mit python3
% nachgerechnet (scripts/lp/trigonometrische-gleichungen/zahlen.py, 08.10.2026, Folgefehler-Fälle eingeschlossen).
\documentclass[11pt]{article}
\usepackage{../lp-druck}
\renewcommand{\lpname}{Trigonometrische Gleichungen}
\renewcommand{\lpteilgebiet}{5.5}
\renewcommand{\lpfuss}{Bewertungspaket}
\newcolumntype{P}{>{\centering\arraybackslash}p{10mm}}
\newenvironment{raster}{\par\smallskip\noindent\tabularx{\linewidth}{@{}>{\raggedright\arraybackslash}p{38mm}P X@{}}\toprule
  \small\textsc{Teilschritt} & \small\textsc{P} & \small\textsc{Musterlösung}\\\midrule}{\bottomrule\endtabularx\par}
\newcommand{\fehler}[1]{\par\smallskip{\small\textcolor{lprot}{\bfseries Typische Fehler:} #1}}
\newcommand{\E}{\,{\footnotesize\bfseries\color{lpgruen}(E)}}
\newcommand{\A}{\,{\footnotesize\bfseries\color{lpblau}(A)}}
\begin{document}
\lpkopf{Bewertungspaket zum Gesamttest}{}

\begin{lpachtung}{Erst lösen, dann öffnen}
Dieses Paket enthält die Musterlösung. Wer es vor dem Lösen liest, misst am Ende nur, wie gut das Abschreiben gelingt.
\end{lpachtung}

\section*{1 · So gehst du vor}
\begin{enumerate}[leftmargin=*,itemsep=2pt]
\item \textbf{Gesamttest lösen} — vollständig, mit Weg und Skizzen, ohne dieses Paket; Teil A ohne, Teil B mit Taschenrechner.
\item \textbf{Fotografieren oder scannen}, gut lesbar, eine Seite pro Bild. \textbf{Kein Name} auf den Bildern.
  Handyfotos tragen oft unsichtbar Standort, Zeit und Gerät mit (Metadaten). Lade darum lieber einen Scan oder ein
  Bildschirmfoto des Fotos hoch, oder entferne vorher die Standortangaben.
\item \textbf{Bewerten lassen.} Ob du dafür eine KI nutzt und welche, entscheidest du selbst und in eigener Verantwortung —
  je nachdem, was dir aufgrund deines Alters und deiner persönlichen Situation erlaubt ist (Altersgrenzen und
  Nutzungsbedingungen des Anbieters, allenfalls Einverständnis deiner Eltern). Lade nur deine Lösung und dieses PDF hoch,
  keine persönlichen Daten, und füge den Auftrag aus Abschnitt~2 ein. Ohne KI kannst du dich mit Abschnitt~3 selbst bewerten.
\item \textbf{Rückmeldung prüfen.} Stimmt, was die KI aus deiner Handschrift und deinen Skizzen gelesen hat? Eine KI kann sich
  irren — im Zweifel gilt das Raster in Abschnitt~3.
\item \textbf{Gezielt wiederholen} nach Abschnitt~4.
\end{enumerate}

\needspace{6\baselineskip}
\section*{2 · Auftrag an die KI}
\begin{tcolorbox}[breakable,colback=lplinie!25,colframe=lplinie,boxrule=0.5pt,arc=1mm,fontupper=\small]
Du bist Korrektorin oder Korrektor für Mathematik an einer Schweizer Berufsmaturitätsschule. Im Anhang findest du (1)~das Bewertungspaket zum Gesamttest «Trigonometrische Gleichungen» mit Musterlösung und Punkteraster und (2)~Fotos meiner handschriftlichen Lösung.\par\medskip
Geh so vor:\par
1. Schreib zu jeder Aufgabe G1 bis G7 zuerst ab, was du in meiner Lösung liest — Weg, Ergebnis und was in der Skizze eingezeichnet ist (welche Gerade, wo die Punkte liegen). Was du nicht sicher lesen kannst, markierst du mit [unsicher]. Nicht raten.\par
2. Vergib die Punkte genau nach dem Raster im Bewertungspaket (Abschnitt 3), ganze Punkte je Zeile. Ziehe einen Fehler nur einmal ab: Rechne ich mit einem falschen Zwischenergebnis richtig weiter, gibt es die Folgepunkte. Ausnahme: Zeilen mit \textbf{(E)} gibt es nur für das richtige Ergebnis, nie als Folgepunkt. Die Zeilen «Typische Fehler» sagen, welche Rasterzeile bei diesem Fehler 0 Punkte gibt — sie sind kein zusätzlicher Abzug.\par
3. Andere richtige Wege zählen voll. Gleichwertige Schreibweisen zählen voll: Dezimalpunkt oder Dezimalkomma; zwischen den Elementen einer Menge der Strichpunkt oder, wenn die Zahlen einen Dezimalpunkt haben, ein Komma (bei Dezimalkomma ist ein Komma als Trenner mehrdeutig — dann zählt, was sich eindeutig lesen lässt, sonst [unsicher]); Winkel im Bogenmass statt in Grad, wenn richtig umgerechnet; $\frac{\sqrt3}{3}$ oder $\frac{1}{\sqrt3}$, $\pm 45^\circ + k \cdot 360^\circ$ statt zweier Formeln, jede andere Grundlösung derselben Familie ($-60^\circ + k \cdot 180^\circ$ statt $120^\circ + k \cdot 180^\circ$), eine Abweichung um $0.1^\circ$ durch Runden — ausser wo das Raster eine bestimmte Form verlangt. In Teil A gilt ein gerundeter Taschenrechnerwert nicht als exakt. Drei Stufen: Zeilen mit \textbf{(A)} verlangen nur Ablesen, diese Punkte gibt es auch ohne Weg. Zeilen mit \textbf{(E)} sind Ergebnispunkte: Das Ergebnis muss stimmen \emph{und} aus einem erkennbaren Weg hervorgehen, ausser die Zeile sagt ausdrücklich etwas anderes. Alle übrigen Zeilen bewerten den Weg selbst.\par
4. Begründe jeden Abzug in einem Satz und nenne die Fehlerart (zum Beispiel «Regel des Sinus beim Cosinus benutzt» oder «Periode des Tangens mit 360° statt 180°»). Schreib die Musterlösung nicht ab — zeig mir, wo mein Fehler liegt.\par
5. Gib zum Schluss eine Tabelle «Aufgabe | Punkte | Begründung», die Gesamtpunktzahl (von 25) und — nach der Selbsteinschätzung in Abschnitt 4 — welches Kapitel des Leitprogramms ich wiederholen soll. Keine Note.\par
6. Fehlt ein Foto oder ist eine Aufgabe unlesbar, sag es, statt sie zu bewerten.
\end{tcolorbox}

\needspace{8\baselineskip}
\section*{3 · Musterlösung und Punkteraster}
\textbf{Allgemein:} Ganze Punkte je Zeile. Ein Fehler wird nur einmal abgezogen (Folgefehler). Gleichwertige Wege und
Schreibweisen zählen voll. In der Spalte P:
\textbf{(A)}~= nur ablesen, zählt auch ohne Weg · \textbf{(E)}~= Ergebnispunkt, Ergebnis richtig
\emph{und} Weg erkennbar (ausser die Zeile sagt ausdrücklich etwas anderes) · ohne Zeichen = die Zeile bewertet den Weg.
Teil A ist ohne, Teil B mit Taschenrechner gestellt; Winkel auf $0.1^\circ$ gerundet. Als Weg zählt auch eine Skizze am
Einheitskreis, die die Gerade und die Lösungspunkte zeigt.
«Typische Fehler» nennt, welche Zeile 0 Punkte gibt, und die Punkte, die bleiben.
\textbf{Teilrichtige Zeilen:} Verlangt eine Zeile mehrere Angaben und stimmt nur ein Teil, gibt die Zeile 0 Punkte —
ausser die Zeile nennt ausdrücklich eine Aufteilung.

\begin{lpaufgabe}{G1}{Am Einheitskreis}{3 P}
\begin{raster}
(a) Gerade & 1 & Die \emph{Senkrechte} $x = \frac{\sqrt3}{2} \approx 0.87$ (der Cosinus ist die $x$-Koordinate). Toleranz: ein halbes Kästchen.\\
(a) Lösungspunkte & 1 & Ihre zwei Schnittpunkte mit dem Kreis, rechts, einer über und einer unter der $x$-Achse, markiert.\\
(b) Lösungsmenge & 1\E & $\cos 30^\circ = \frac{\sqrt3}{2}$, Spiegelbild an der $x$-Achse $360^\circ - 30^\circ$: $\mathbb{L} = \{30^\circ;\ 330^\circ\}$. Beide Winkel nötig.\\
\end{raster}
\fehler{Waagrechte $y = \frac{\sqrt3}{2}$ gezeichnet (Sinus statt Cosinus): beide (a)-Zeilen 0; (b) $\{60^\circ;\ 120^\circ\}$ ist dann ebenfalls 0 (E).
(b) $\{30^\circ;\ 150^\circ\}$ (Regel des Sinus): Zeile (b) 0.
(b) $\{-30^\circ;\ 30^\circ\}$: Zeile (b) 0, denn $-30^\circ$ liegt nicht im Intervall.}
\end{lpaufgabe}

\begin{lpaufgabe}{G2}{Wie viele Lösungen?}{3 P}
\begin{raster}
(a) & 1 & Keine: Der Sinus liegt zwischen $-1$ und $1$; die Waagrechte $y = -1.2$ verfehlt den Kreis. $\mathbb{L} = \{\,\}$. Zahl \emph{und} Begründung nötig.\\
(b) & 1 & Zwei: Die Senkrechte $x = -1$ berührt den Kreis nur in $(-1 \mid 0)$, bei $180^\circ$ — eine Lösung je Umdrehung. In $[0^\circ;\, 720^\circ[$ sind es zwei Umdrehungen: $180^\circ$ und $540^\circ$. Zahl und Begründung nötig.\\
(c) & 1 & Zwei: Der Punkt $S(1 \mid -50)$ existiert, die Gerade durch $O$ und $S$ trifft den Kreis in zwei Gegenpunkten (bei rund $91.1^\circ$ und $271.1^\circ$ — die Winkel sind nicht verlangt). Zahl und Begründung nötig.\\
\end{raster}
\fehler{(b) «eine, $180^\circ$» (zweite Umdrehung übersehen): Zeile (b) 0. (b) «zwei, weil $-1$ zwischen $-1$ und $1$ liegt» (Berührung als Schnitt gezählt, Begründung falsch): Zeile (b) 0. (c) «keine, weil $-50 < -1$» (Grenze des Sinus auf den Tangens übertragen): Zeile (c) 0. Richtige Zahl ohne Begründung: die Zeile 0.}
\end{lpaufgabe}

\begin{lpaufgabe}{G3}{Alle Lösungen}{4 P}
\begin{raster}
(a) Grundlösung & 1\E & $\tan 60^\circ = \sqrt3$; negativ im II.~Quadranten: $180^\circ - 60^\circ = 120^\circ$ (oder $-60^\circ$, oder $300^\circ$).\\
(a) Periode & 1 & $\varphi = 120^\circ + k \cdot 180^\circ$, $k \in \mathbb{Z}$. Folgepunkt: gilt mit einer falschen Grundlösung, wenn $+\,k \cdot 180^\circ$ steht.\\
(b) Grundlösungen & 1\E & $\cos 45^\circ = \frac{\sqrt2}{2}$ und $360^\circ - 45^\circ = 315^\circ$ (oder $\pm 45^\circ$). Beide nötig.\\
(b) Periode & 1 & $\varphi = 45^\circ + k \cdot 360^\circ$ oder $\varphi = 315^\circ + k \cdot 360^\circ$, $k \in \mathbb{Z}$ (gleichwertig $\varphi = \pm 45^\circ + k \cdot 360^\circ$). Folgepunkt bei falschen oder fehlenden Grundlösungen, wenn jede angegebene Familie $+\,k \cdot 360^\circ$ trägt.\\
\end{raster}
\fehler{(a) $\varphi = 120^\circ + k \cdot 360^\circ$ (Periode des Sinus): Periodenzeile 0 $\to$ 1 von 2\,P. Wer zwei Familien $120^\circ + k \cdot 360^\circ$ und $300^\circ + k \cdot 360^\circ$ schreibt, hat alle Lösungen: volle Punkte.
(a) Grundlösung $60^\circ$ (Vorzeichen übersehen): Zeile Grundlösung 0, Periode als Folgepunkt 1 $\to$ 1 von 2\,P.
(b) nur $45^\circ + k \cdot 360^\circ$ (zweite Grundlösung fehlt): Grundlösungen-Zeile 0, Periode 1 $\to$ 1 von 2\,P. (b) $+\,k \cdot 180^\circ$: Periodenzeile 0.
Fehlt «$k \in \mathbb{Z}$» ganz, aber ist $k$ als ganze Zahl gemeint: kein Abzug. Im Bogenmass richtig ($-\frac{\pi}{3} + k\pi$; $\pm\frac{\pi}{4} + 2k\pi$): kein Abzug.}
\end{lpaufgabe}

\begin{lpaufgabe}{G4}{Sinusgleichung in einem anderen Intervall}{4 P}
\begin{raster}
Hauptwert & 1 & $\varphi_1 = \sin^{-1}(-0.35) \approx -20.5^\circ$ — liegt schon in $[-180^\circ;\, 180^\circ[$.\\
zweite Lösung & 1\E & Spiegelung an der $y$-Achse: $180^\circ - (-20.5^\circ) = 200.5^\circ$ (gleichwertig: direkt $-180^\circ + 20.5^\circ = -159.5^\circ$).\\
Lösungsmenge & 1 & $200.5^\circ - 360^\circ = -159.5^\circ$; $\mathbb{L} = \{-159.5^\circ;\ -20.5^\circ\}$, aufsteigend, beide im Intervall. Folgepunkt: gilt mit falschen Werten, wenn die eigenen zwei Werte richtig ins Intervall gebracht und aufsteigend geordnet sind. Fehlt eine Lösung, gibt die Zeile 0.\\
Probe & 1 & Beide Punkte liegen unter der $x$-Achse (III. und IV.~Quadrant), der Sinus ist dort negativ — als Skizze im Kreis oder als Satz. Folgepunkt: gilt, wenn die Probe an den eigenen Werten richtig ausgeführt ist, auch wenn sie dabei einen Fehler aufdeckt.\\
\end{raster}
\fehler{$\mathbb{L} = \{-20.5^\circ;\ 200.5^\circ\}$ ($200.5^\circ$ nicht ins Intervall gebracht) oder $\{200.5^\circ;\ 339.5^\circ\}$ (Intervall $[0^\circ;\, 360^\circ[$ statt des verlangten): Zeile «Lösungsmenge» 0 $\to$ 3 von 4\,P.
Nur $-20.5^\circ$ (zweite Lösung fehlt, nur den Wert des Rechners angegeben): «zweite Lösung» 0 und «Lösungsmenge» 0 $\to$ 2 von 4\,P mit Probe, sonst 1 von 4\,P.
$360^\circ - (-20.5^\circ) = 380.5^\circ$, ins Intervall $20.5^\circ$ (Regel des Cosinus): «zweite Lösung» 0 (E); «Lösungsmenge» 1 als Folgepunkt; «Probe» 1 nur, wenn sie benennt, dass $20.5^\circ$ über der $x$-Achse liegt $\to$ 3 bzw. 2 von 4\,P.
Vorzeichen übersehen ($\sin^{-1}(0.35)$, $\{20.5^\circ;\ 159.5^\circ\}$): «Hauptwert» 0; «zweite Lösung» 0, weil (E) nur das richtige Ergebnis zählt (kein zusätzlicher Abzug, sondern die Regel für (E)); «Lösungsmenge» 1 als Folgepunkt; «Probe» 1 nur, wenn sie den Widerspruch benennt (Punkte über der $x$-Achse, Sinus positiv) $\to$ 2 bzw. 1 von 4\,P.}
\end{lpaufgabe}

\begin{lpaufgabe}{G5}{Alle Lösungen, dann ein Intervall}{4 P}
\begin{raster}
Hauptwert & 1 & $\varphi_1 = \cos^{-1}(0.42) \approx 65.2^\circ$\\
zweite Lösung & 1\E & Spiegelung an der $x$-Achse: $360^\circ - 65.2^\circ = 294.8^\circ$ (oder $-65.2^\circ$).\\
(a) Periode & 1 & $\varphi = 65.2^\circ + k \cdot 360^\circ$ oder $\varphi = 294.8^\circ + k \cdot 360^\circ$, $k \in \mathbb{Z}$ (gleichwertig $\varphi = \pm 65.2^\circ + k \cdot 360^\circ$). Folgepunkt: gilt mit falschen oder fehlenden Grundlösungen, wenn jede angegebene Familie $+\,k \cdot 360^\circ$ trägt.\\
(b) Lösungsmenge & 1 & $k = -1$ und $k = 0$: $\mathbb{L} = \{-294.8^\circ;\ -65.2^\circ;\ 65.2^\circ;\ 294.8^\circ\}$, aufsteigend, nichts ausserhalb von $[-360^\circ;\, 360^\circ[$. Folgepunkt: gilt, wenn die Menge genau die Werte der eigenen Familien im Intervall enthält, aufsteigend. Fehlt eine Lösung, gibt die Zeile 0.\\
\end{raster}
\fehler{$180^\circ - 65.2^\circ = 114.8^\circ$ (Regel des Sinus): «zweite Lösung» 0; mit $114.8^\circ$ richtig weitergerechnet ($\{-294.8^\circ;\ -245.2^\circ;\ 65.2^\circ;\ 114.8^\circ\}$) gelten die Folgepunkte $\to$ 3 von 4\,P.
Nur $\varphi = 65.2^\circ + k \cdot 360^\circ$ (zweite Lösung fehlt): «zweite Lösung» 0, «Periode» 1, «Lösungsmenge» 0 $\to$ 2 von 4\,P.
(b) nur $\{65.2^\circ;\ 294.8^\circ\}$ oder nur $\{-294.8^\circ;\ -65.2^\circ\}$ (eine Umdrehung des Intervalls vergessen): «Lösungsmenge» 0 $\to$ 3 von 4\,P.
(a) $+\,k \cdot 180^\circ$: «Periode» 0; (b) gilt als Folgepunkt nur, wenn die Menge zu diesen Familien passt.}
\end{lpaufgabe}

\begin{lpaufgabe}{G6}{Tangens über zwei Umdrehungen}{3 P}
\begin{raster}
Hauptwert & 1 & $\varphi_1 = \tan^{-1}(-2.4) \approx -67.4^\circ$ — liegt nicht im Intervall.\\
erste Umdrehung & 1\E & $-67.4^\circ + 180^\circ = 112.6^\circ$ und $-67.4^\circ + 360^\circ = 292.6^\circ$. Beide nötig.\\
Lösungsmenge & 1 & Periode $180^\circ$ weiter: $472.6^\circ$, $652.6^\circ$. $\mathbb{L} = \{112.6^\circ;\ 292.6^\circ;\ 472.6^\circ;\ 652.6^\circ\}$, aufsteigend, nichts ausserhalb von $[0^\circ;\, 720^\circ[$. Folgepunkt: gilt, wenn die Menge aus den eigenen Werten mit Schritt $180^\circ$ richtig gebildet ist. Fehlt eine Lösung, gibt die Zeile 0.\\
\end{raster}
\fehler{Periode $360^\circ$ gedacht ($\{112.6^\circ;\ 472.6^\circ\}$ oder $\{292.6^\circ;\ 652.6^\circ\}$): «erste Umdrehung» 0 und «Lösungsmenge» 0 $\to$ 1 von 3\,P.
Nur $\{112.6^\circ;\ 292.6^\circ\}$ (zweite Umdrehung vergessen): «Lösungsmenge» 0 $\to$ 2 von 3\,P.
$-67.4^\circ$ in der Menge (Hauptwert stehen gelassen): «Lösungsmenge» 0 $\to$ 2 von 3\,P.
$180^\circ - (-67.4^\circ) = 247.4^\circ$ (Regel des Sinus): «erste Umdrehung» 0, «Lösungsmenge» 0 $\to$ 1 von 3\,P.}
\end{lpaufgabe}

\begin{lpaufgabe}{G7}{Fehler finden}{4 P}
\begin{raster}
(a) Fehler benennen & 1 & Der Schritt $\varphi_2 = 180^\circ - \varphi_1$: Das ist die Regel des Sinus.\\
(a) Begründung & 1 & Bei $53.1^\circ$ liegt der Punkt rechts der $y$-Achse, der Cosinus ist positiv ($\cos 53.1^\circ \approx 0.6$), gesucht ist $-0.6$. Beim Cosinus liegen die zwei Punkte auf derselben Senkrechten $x = -0.6$, gespiegelt an der $x$-Achse.\\
(b) zweite Lösung & 1\E & $\varphi_2 = 360^\circ - 126.9^\circ = 233.1^\circ$\\
(b) Lösungsmenge & 1 & $\mathbb{L} = \{126.9^\circ;\ 233.1^\circ\}$, aufsteigend. Folgepunkt: passt die Menge zur eigenen zweiten Lösung.\\
\end{raster}
\fehler{«$126.9^\circ$ ist falsch» (der Wert des Rechners stimmt): Zeile (a) Fehler 0.
Begründung nur «die Regel ist falsch» ohne Bezug zum Kreis oder zum Vorzeichen: Begründungs-Zeile 0.
(b) $180^\circ + 126.9^\circ = 306.9^\circ$: Zeile «zweite Lösung» 0 ($\cos 306.9^\circ > 0$).}
\end{lpaufgabe}

\needspace{14\baselineskip}
\section*{4 · Selbsteinschätzung}
\begin{tabularx}{\linewidth}{@{}l X@{}}\toprule
22 – 25 P & Die geprüften Teile sitzen. Wo du Punkte verloren hast: das Kapitel dieser Aufgabe nochmals (Zuordnung unten).\\
17 – 21 P & Den schwächsten Teil nochmals: Tüfteln und Übungen des Kapitels, in dem du die meisten Punkte verloren hast.\\
11 – 16 P & Zurück zu den Kapiteln aller Aufgaben, in denen du Punkte verloren hast.\\
0 – 10 P & Zurück zu Kapitel 1 und von dort der Reihe nach weiter.\\\midrule
\multicolumn{2}{@{}p{\linewidth}@{}}{\small\textbf{Aufgabe $\to$ Kapitel:} G1 $\to$ Kapitel 1; G2 $\to$ Kapitel 1, 3, 4; G3 $\to$ Kapitel 1, 3, 4; G4 $\to$ Kapitel 2, 4; G5 $\to$ Kapitel 2, 4; G6 $\to$ Kapitel 3, 4; G7 $\to$ Kapitel 2}\\\bottomrule
\end{tabularx}
\end{document}
